\SourceLeaf{174}Princips wegen nicht für gerathen halten, wenn ich auch unter den jetzigen günstigen Umständen gern die Hoffnung theile, dass das Minus von 3448 Thalern schon in den nächsten Jahren sich bedeutend vermindern würde.

\medskip
Ich glaube im Vorstehenden der Universitäts-Kirchendeputation das hinlängliche Material zusammengebracht zu haben, wonach mit bewusster Sicherheit ein Beschluss gefasst werden kann, bin aber gern zu weitern Erläuterungen bereit, wenn solche für nöthig gehalten werden sollten.

Göttingen den 19. October 1851.\hfill C. F. \textsc{Gauss}.
\begin{center}\rule{.28\linewidth}{.3pt}\end{center}

\Needspace{75pt}
\pdfbookmark[1]{Berechnung der Mittelwerthe der Beiträge und Pensionen}{witwen-mittelwerthe}
\begin{center}[\textit{Berechnung der Mittelwerthe der Beiträge und Pensionen.}\end{center}
\noindent\hspace*{\parindent}Die Anzahl der in ununterbrochener Reihenfolge nach der Numerirung ihrer Beitrittszeit bis zum 1. October 1851 verstorbenen oder auf andere Weise aus der Professoren-Witwenkasse ausgeschiedenen Mitglieder beträgt 134. Für jedes dieser 134 ausgeschiedenen Mitglieder ist in folgender Tabelle zusammengestellt: das Jahr an dessen 1. Oct. (und nur ausnahmsweise an dessen 1. Apr. bei Nr. 134) der Beitritt erfolgte, die Anzahl der Jahre und Monate, während welcher die Beiträge oder die Witwen und Waisenpensionen ein- oder ausgezahlt wurden, ebenso diejenigen welche vergingen bis zu dem Zeitpunkte von wo an die Pension gerechnet wird. Aus diesen Daten sind die in den daneben stehenden Spalten enthaltenen Werthe bestimmt, welche im Zeitpunkte des Eintritts für den Zinsfuss von \(3\tfrac{1}{2}\) Proc. und 4 Proc. den jährlich mit 10 Thl. eingezahlten Beiträgen und den mit 250 Thl. ausgezahlten Witwen- und Waisenpensionen gleichkommen. Die Ansätze zu 10 Thl. und 250 Thl. sind von \textsc{Gauss} wohl deshalb gewählt, weil hierfür schon der grösste Theil der Tafel im Jahre 1845 berechnet war. Für die am 1. Oct. 1851 noch lebenden Witwen der Mitglieder Nr. 86. 103. 104. 109. 112. 114. 122 und 134 ist der Capitalwerth derjenigen Pensionsbezüge die nach jenem Zeitpunkte noch erfolgen mussten, mit Zuhülfenahme der \textsc{Brune}’schen Sterblichkeitstafeln bestimmt und unter der Voraussetzung dass die Zahlung am Ende jedes Jahres aber nur bis zum Schluss des Sterbemonats erfolgt. Die aus dieser Zusammenstellung sich unmittelbar ergebenden mittleren Capitalwerthe der Beiträge und Pensionen sind in der Bilanzrechnung von 1851 benutzt und zwar nachdem durch Multiplication mit einem Correctionsfactor der Umstand berücksichtigt ist, dass die Pensionen nicht jährlich sondern halbjährlich gezahlt werden. Für die Waisenpensionen ist hier angenommen, dass sie bis zum Schlusse des Monats galten, in welchem das 20\textsuperscript{ste} Lebensjahr des Kindes vollendet wurde.

Die handschriftliche Tafel enthält keine Überschrift der einzelnen Spalten wie hier der Abdruck.]
